% ---
% title: Vandermonde's identity
% tags: combinatorial-proof, combinations
% ---

\begin{problem}
	Prove Vandermonde's identity
	\[ \sum_{j=0}^k\binom nj\binom m{k-j}=\binom{n+m}k \]
	by a counting argument (apples and pears).
\end{problem}

\begin{solution}
	\idea Count the same set in two ways, as in Exercise~\exref{pascals-rule}.
	\step{Setting.} A bowl contains $n$ different apples and $m$ different
	pears. Count the ways to choose $k$ fruits.
	\step{Directly.} $k$ out of $n+m$ fruits: $\binom{n+m}k$.
	\step{By the number of apples.} Let $j$ be the number of apples chosen
	($j=0,1,\dots,k$). Then $k-j$ pears are chosen: $\binom nj\binom m{k-j}$
	ways. The cases $j=0,\dots,k$ are disjoint and cover all choices, so the
	total is
	\[ \sum_{j=0}^k\binom nj\binom m{k-j}. \]
	\begin{answer}
		Both expressions count the $k$-element selections from $n+m$ fruits,
		hence
		\[ \sum_{j=0}^k\binom nj\binom m{k-j}=\binom{n+m}k. \]
	\end{answer}
	\begin{remark}
		Terms with $j>n$ or $k-j>m$ are $0$ by convention, so the identity holds
		for all $k$. Check with $n=4$, $m=5$, $k=3$:
		\[ 1\cdot10+4\cdot10+6\cdot5+4\cdot1=84=\binom93. \]
		Pascal's rule (Exercise~\exref{pascals-rule}) is the special case $m=1$.
	\end{remark}
\end{solution}
